\documentclass[11pt]{article}

\usepackage[margin=1in]{geometry}
\usepackage{amsmath, amssymb, amsthm}
\usepackage{graphicx}
\usepackage{hyperref}
\usepackage{enumitem}

\title{Feature-wise Hyperdimensional Encoding with Regression-based MCI for Data Explanation}
\author{Yuchen}
\date{April 2026}

\begin{document}

\maketitle




\section{Feature-wise HDC Encoding}

Given input $x = (x_1, \dots, x_d)$, we encode each feature independently:

\begin{equation}
z_j(x_j) = \sqrt{\frac{2}{D_j}} \cos(W_j x_j + b_j)
\end{equation}

The full representation is:

\begin{equation}
z(x) = [z_1(x_1), \dots, z_d(x_d)]
\end{equation}

Feature-wise HDC encoder may help with calculating $w_S^*$ in \ref{eq:regression_solution} efficiently. 

\section{Regression Model}

We use ridge regression:

\begin{equation}
\hat{y} = w^\top z(x)
\end{equation}

For a subset $S \subseteq \{1,\dots,d\}$:

\begin{equation}
z_S(x) = [z_j(x_j)]_{j \in S}
\end{equation}

\begin{equation}
w_S^* = \arg\min_w \sum_{n=1}^N (y_n - w^\top z_S(x_n))^2 + \lambda \|w\|^2
\end{equation}

Closed-form solution:

\begin{equation}
w_S^* = (Z_S^\top Z_S + \lambda I)^{-1} Z_S^\top y
\label{eq:regression_solution}
\end{equation}

We don't need to retrain the optimal model given subset $S$ anymore. 

\section{Predictive Power Estimation}

We define predictive power of features in subset $S$ as:

\begin{equation}
\nu(S) = \mathrm{Var}(y) - \mathrm{MSE}(S)
\end{equation}

where:

\begin{equation}
\mathrm{MSE}(S) = \frac{1}{N} \sum_{n=1}^N (y_n - \hat{y}_n^{(S)})^2
\end{equation}

\section{Marginal Contribution Importance (MCI)}

MCI is approximated as:

\begin{equation}
\widehat{I}(i) = \max_{S \subseteq F \setminus \{i\}} \left[\nu(S \cup \{i\}) - \nu(S)\right]
\end{equation}


We highlight a key property of marginal contribution importance (MCI): it provides an upper bound on predictive power.


For any subset of features $S \subseteq F$, let $\nu(S)$ denote predictive power and define normalized MCI as:
\begin{equation}
\tilde I(i) = \frac{I_\nu(i)}{\nu(F)}.
\end{equation}

Then for any subset $S \subseteq F$, we have:
\begin{equation}
\frac{\nu(S)}{\nu(F)} \le \sum_{i \in S} \tilde I(i).
\end{equation}

The left-hand side represents the fraction of total predictive power explained by $S$, 
while the right-hand side provides an upper bound based on the maximal marginal contributions 
of individual features. 


% In practice, we estimate this via sampled permutations:

% \begin{equation}
% \widehat{I}(i) = \max_{\sigma} \left[\nu(S_i^\sigma \cup \{i\}) - \nu(S_i^\sigma)\right]
% \end{equation}

\section{Computational Advantages}

The feature-wise structure enables:

\begin{itemize}
    \item Block-wise design matrices $Z_S$
    \item Efficient ridge updates
    \item Avoidance of full retraining
\end{itemize}

\section{Discussion}

Our framework provides:

\begin{itemize}
    \item Nonlinear modeling via RFF
    \item Feature modularity via HDC
    \item Efficient MCI approximation via regression
\end{itemize}


\bibliographystyle{plain}
\bibliography{references}

\end{document}